\chapter*{Appendices}
\addcontentsline{toc}{chapter}{Appendices}
\markboth{APPENDICES}{APPENDICES}

\section*{Appendix A: Operational Installation and Deployment Guide}
To facilitate the deployment, replication, and testing of the \textbf{Predictra Threat Map} platform, this appendix provides a complete step-by-step developer installation guide.

\subsection*{A.1 Prerequisites and Runtime Environments}
Ensure the target deployment virtual machine or local workstation is equipped with the following dependencies:
\begin{itemize}
    \item \textbf{Node.js Runtime:} Version v20.x.x LTS or higher.
    \item \textbf{Package Manager:} npm version 10.x.x or higher.
    \item \textbf{Database:} MongoDB Atlas cloud database cluster (M1 tier or higher) or a local MongoDB community edition service.
\end{itemize}

\subsection*{A.2 Backend Configuration Variables}
Create a configuration file named `.env` in the `/backend` root directory. The file must contain the following variables:
\begin{verbatim}
PORT=3001
MONGO_URI=mongodb://<user>:<pwd>@cluster.mongodb.net/db
NODE_ENV=production
BATCH_INTERVAL_MS=300
MAX_PENDING=500
BATCH_DB_INSERT=25
\end{verbatim}

\subsection*{A.3 Installation and Local Execution Commands}
To launch the complete platform locally in development mode:
\begin{enumerate}
    \item \textbf{Clone the Repository:}
    \begin{verbatim}
    git clone https://github.com/predictra/predictra-threat-map.git
    cd predictra-threat-map
    \end{verbatim}
    \item \textbf{Install Root and Package Dependencies:}
    \begin{verbatim}
    npm install
    \end{verbatim}
    \item \textbf{Launch Backend Server (Port 3001):}
    \begin{verbatim}
    cd backend
    npm install
    npm run dev
    \end{verbatim}
    \item \textbf{Launch Frontend App (Port 5173):}
    \begin{verbatim}
    cd ../frontend
    npm install
    npm run dev
    \end{verbatim}
\end{enumerate}

\subsection*{A.4 Production Process Management (PM2)}
In production environments, the backend server must be managed using the PM2 process manager to ensure zero-downtime clustering and automated crash recoveries:
\begin{verbatim}
pm2 start server.js --name "predictra-backend" -i max --watch
pm2 save
pm2 startup
\end{verbatim}

---

\section*{Appendix B: API Integrations and Data Streams Reference}
This section provides a quick-reference catalog of the internal REST API endpoints exposed by the Node.js Express server to facilitate security tool integrations (e.g. connecting the Threat Map feeds to internal corporate SIEM dashboards).

\subsection*{B.1 Threat Feeds Stream Endpoint}
\begin{itemize}
    \item \textbf{Endpoint:} \texttt{GET /api/feed}
    \item \textbf{Protocol:} Server-Sent Events (SSE)
    \item \textbf{Response Headers:}
    \begin{verbatim}
    Content-Type: text/event-stream
    Cache-Control: no-cache
    Connection: keep-alive
    \end{verbatim}
    \item \textbf{Event Type:} \texttt{event: attacks}
    \item \textbf{Payload Example:}
    \begin{verbatim}
    data: [
      {
        "a_c": 1,
        "a_n": "Cobalt Strike C2 Active",
        "a_t": "malware",
        "s_ip": "185.220.101.5",
        "s_co": "DE",
        "s_la": 51.2993,
        "s_lo": 9.491,
        "d_ip": "192.168.10.4",
        "d_co": "US",
        "d_la": 37.0902,
        "d_lo": -95.7129,
        "source_api": "c2tracker",
        "timestamp": "2026-05-21T12:00:00.000Z"
      }
    ]
    \end{verbatim}
\end{itemize}

\subsection*{B.2 Historical Intelligence Queries API}
\begin{itemize}
    \item \textbf{Endpoint:} \texttt{GET /api/history}
    \item \textbf{Query Parameters:}
    \begin{itemize}
        \item \texttt{page} (Number): Page index (default: \texttt{1})
        \item \texttt{limit} (Number): Page size limit (default: \texttt{50})
        \item \texttt{query} (String): Search keyword matching country, sector, or IP
    \end{itemize}
    \item \textbf{Response Format:} Paginated JSON array containing BSON documents matching database schemas.
\end{itemize}

---

\section*{Appendix C: Automated Integration Test Suites (Jest \& Supertest)}
To guarantee backend regression safety, route integrity, and database transaction safety, the Predictra platform implements automated test suites using the \textbf{Jest} unit testing framework and the \textbf{Supertest} HTTP assertion library.

\subsection*{C.1 Ingestion API Integration Test Config (\texttt{tests/api.test.js})}
The backend API testing environment mocks active database calls using an in-memory MongoDB service (\texttt{mongodb-memory-server}), checking that geocoding routes and threat storage pipelines operate correctly without writing fake logs to production collections:
\begin{verbatim}
const request = require('supertest');
const mongoose = require('mongoose');
const { MongoMemoryServer } = require('mongodb-memory-server');
const app = require('../server_app'); // Express App Core
const ThreatEvent = require('../models/ThreatEvent');

let mongoServer;

beforeAll(async () => {
  mongoServer = await MongoMemoryServer.create();
  const uri = mongoServer.getUri();
  await mongoose.disconnect(); // Avoid active database conflicts
  await mongoose.connect(uri, {
    useNewUrlParser: true,
    useUnifiedTopology: true
  });
});

afterAll(async () => {
  await mongoose.disconnect();
  await mongoServer.stop();
});

beforeEach(async () => {
  await ThreatEvent.deleteMany({}); // Flush collections between tests
});

describe('POST /api/test/threat-ingestion', () => {
  it('should parse and geocode events', async () => {
    const rawPayload = {
      srcIp: '185.220.101.5',
      dstIp: '8.8.8.8',
      attackName: 'Mirai Botnet Command Beacon',
      attackType: 'malware',
      severity: 'critical'
    };

    const response = await request(app)
      .post('/api/test/threat-ingestion')
      .send(rawPayload)
      .expect(201);

    expect(response.body).toHaveProperty('success', true);
    expect(response.body).toHaveProperty('data');
    expect(response.body.data.a_n).toBe('Mirai Botnet Command Beacon');

    // Confirm collection persistence
    const savedEvent = await ThreatEvent
      .findOne({ s_ip: '185.220.101.5' });
    expect(savedEvent).toBeDefined();
    expect(savedEvent.a_t).toBe('malware');
    expect(savedEvent.s_co).toBe('DE'); // Geocoder translation check
  });

  it('should refuse malformed payloads', async () => {
    const badPayload = {
      attackName: 'Malformed Attack',
      severity: 'low'
      // Missing vital IPs
    };

    await request(app)
      .post('/api/test/threat-ingestion')
      .send(badPayload)
      .expect(400);

    const count = await ThreatEvent.countDocuments();
    expect(count).toBe(0);
  });
});
\end{verbatim}

---

\section*{Appendix D: Vercel CDN and Reverse-Proxy Configurations}
To orchestrate secure client-to-server communications and bypass standard browser
CORS (\textbf{Cross-Origin Resource Sharing}) restrictions without compromising
security headers, we implemented a Vercel reverse-proxy configuration.

\subsection*{D.1 Production Vercel Router (\texttt{frontend/vercel.json})}
The production file located in the frontend root ensures that all visual components issue requests directly to the same CDN origin host, routing backend queries behind internal proxy links to prevent CORS failures:
\begin{verbatim}
{
  "version": 2,
  "rewrites": [
    {
      "source": "/api/feed",
      "destination": "https://api.predictra.io/api/feed"
    },
    {
      "source": "/api/history",
      "destination": "https://api.predictra.io/api/history"
    },
    {
      "source": "/api/(.*)",
      "destination": "https://api.predictra.io/api/$1"
    },
    {
      "source": "/(.*)",
      "destination": "/index.html"
    }
  ],
  "headers": [
    {
      "source": "/(.*)",
      "headers": [
        {
          "key": "X-Content-Type-Options",
          "value": "nosniff"
        },
        {
          "key": "X-Frame-Options",
          "value": "DENY"
        },
        {
          "key": "X-XSS-Protection",
          "value": "1; mode=block"
        },
        {
          "key": "Referrer-Policy",
          "value": "strict-origin-when-cross-origin"
        }
      ]
    }
  ]
}
\end{verbatim}

---

\section*{Appendix E: Global Frontend Zustand Storage Engine}
To maintain a high framerate on our 3D visualization canvas while receiving dense threat updates over SSE channels, the client application relies on a decoupled, centralized Zustand state manager.

\subsection*{E.1 Core Stream Store (\texttt{stream/useStreamStore.ts})}
The client store acts as the central state hub, handling server-sent events, managing active canvas arc counts, tracking rendering frame rates, and dropping incoming threat events if performance falls below acceptable thresholds:
\begin{verbatim}
import { create } from 'zustand';
import { RingBuffer } from '../utils/perf';

interface ThreatEvent {
  id: string;
  a_n: string;
  a_t: string;
  s_co: string;
  s_la: number;
  s_lo: number;
  d_co: string;
  d_la: number;
  d_lo: number;
  source_api: string;
  timestamp: string;
}

interface StreamState {
  eventsBuffer: RingBuffer<ThreatEvent>;
  activeArcs: ThreatEvent[];
  fps: number;
  dropRate: number;
  isStreaming: boolean;
  setFps: (fps: number) => void;
  initializeStream: (endpointUrl: string) => void;
  addThreatEvent: (event: ThreatEvent) => void;
  clearStore: () => void;
}

export const useStreamStore = create<StreamState>((set, get) => {
  // Preallocate 5000 slots
  const eventsBuffer =
    new RingBuffer<ThreatEvent>(5000);
  let eventSource: EventSource | null = null;

  return {
    eventsBuffer,
    activeArcs: [],
    fps: 60,
    dropRate: 0.0,
    isStreaming: false,

    setFps: (fps) => {
      let dropRate = 0.0;
      if (fps < 30) {
        dropRate = 0.75; // Heavy load: drop 75% visual arcs
      } else if (fps < 48) {
        dropRate = 0.40; // Moderate load: drop 40% visual arcs
      }
      set({ fps, dropRate });
    },

    initializeStream: (endpointUrl) => {
      if (eventSource) return;

      set({ isStreaming: true });
      eventSource = new EventSource(endpointUrl);

      eventSource.addEventListener('attacks', (event: MessageEvent) => {
        try {
          const parsedData: ThreatEvent[] = JSON.parse(event.data);
          parsedData.forEach((threat) => {
            get().addThreatEvent(threat);
          });
        } catch (err) {
          console.error('Error parsing streaming event:', err);
        }
      });

      eventSource.onerror = (err) => {
        console.error('EventSource connection error:', err);
        eventSource?.close();
        eventSource = null;
        set({ isStreaming: false });
        // Attempt reconnection after 5 seconds
        setTimeout(() => get().initializeStream(endpointUrl), 5000);
      };
    },

    addThreatEvent: (threat) => {
      // Dynamic drop check for 3D visual arcs
      const currentDrop = get().dropRate;
      const shouldRenderVisual = Math.random() >= currentDrop;

      eventsBuffer.push(threat);

      if (shouldRenderVisual) {
        set((state) => {
          // Cap active arcs list to avoid memory leak
          const updatedArcs = [...state.activeArcs, threat];
          if (updatedArcs.length > 150) {
            updatedArcs.shift(); // Remove oldest arc line
          }
          return { activeArcs: updatedArcs };
        });
      }
    },

    clearStore: () => {
      if (eventSource) {
        eventSource.close();
        eventSource = null;
      }
      eventsBuffer.clear();
      set({ activeArcs: [], isStreaming: false });
    }
  };
});
\end{verbatim}

---

\section*{Appendix F: Sample Threat Event (BSON vs Enriched JSON)}
To illustrate the data transformation performed by our backend enrichment pipeline, this section contrasts the highly compact, space-saving BSON document stored in the MongoDB Atlas database with the fully resolved, logically structured JSON document exposed to the client application.

\subsection*{F.1 Compact BSON Document Database Structure (82 Bytes)}
This optimized model reduces the storage footprint of each record by over $55\%$, enabling millions of threats to be cached within the database without exceeding storage boundaries:
\begin{verbatim}
{
  "_id": {"$oid": "664c8d1bf8e10b1a084cf4e7"},
  "a_c": 1,
  "a_n": "LockBit Ransomware Activity",
  "a_t": "malware",
  "s_ip": "103.24.15.12",
  "s_co": "CN",
  "s_la": 39.9042,
  "s_lo": 116.4074,
  "d_ip": "194.26.135.8",
  "d_co": "RU",
  "d_la": 55.7558,
  "d_lo": 37.6173,
  "source_api": "ransomwatch",
  "timestamp": {"$date": "2026-05-21T12:00:00.000Z"}
}
\end{verbatim}

\subsection*{F.2 Enriched Logical Client JSON Format (246 Bytes)}
When fetched by the frontend via the historical REST API or stream channels, the controller dynamically expands the document, geolocating the addresses, mapping coordinates, and resolving the target industry sectors:
\begin{verbatim}
{
  "id": "664c8d1bf8e10b1a084cf4e7",
  "attackCount": 1,
  "threatName": "LockBit Ransomware Activity",
  "threatType": "malware",
  "source": {
    "ip": "103.24.15.12",
    "countryCode": "CN",
    "countryName": "China",
    "coordinates": {
      "latitude": 39.9042,
      "longitude": 116.4074
    },
    "organization": "China Telecom Backbone Service"
  },
  "destination": {
    "ip": "194.26.135.8",
    "countryCode": "RU",
    "countryName": "Russian Federation",
    "coordinates": {
      "latitude": 55.7558,
      "longitude": 37.6173
    },
    "organization": "Moscow Financial Services Holding",
    "industrySector": "Financial Services"
  },
  "enrichmentMetadata": {
    "classificationLayer": 2,
    "confidenceScore": 0.95,
    "rdapResolved": true,
    "rdapCacheAgeSeconds": 1420
  },
  "sourceApiFeed": "ransomwatch",
  "timestamp": "2026-05-21T12:00:00.000Z"
}
\end{verbatim}

---

\section*{Appendix G: Geocoding and Coordinate Mappings}
To facilitate real-time mathematical operations in the client's web browser, the platform implements localized 3D coordinate translations, spherical linear interpolation, and geocoding transformations. This appendix provides the complete TypeScript source code for the geocoding and mapping utility.

\subsection*{G.1 Geocoding and 3D Projection (\texttt{utils/geo.ts})}
The utility below translates spherical coordinates into Three.js 3D space, computes great-circle intermediate nodes dynamically using parabolic elevation offsets, and maps standard country names to ISO Alpha-2 codes:
\begin{verbatim}
import * as THREE from 'three';

const GLOBE_RADIUS = 1.0;
const DEG2RAD = Math.PI / 180;

/**
 * Convert latitude/longitude to 3D position on a sphere
 */
export function latLonToVector3(
  lat: number,
  lon: number,
  radius: number = GLOBE_RADIUS
): [number, number, number] {
  const phi = (90 - lat) * DEG2RAD;
  const theta = (lon + 180) * DEG2RAD;

  const x = -(radius * Math.sin(phi) * Math.cos(theta));
  const y = radius * Math.cos(phi);
  const z = radius * Math.sin(phi) * Math.sin(theta);

  return [x, y, z];
}

// Reusable vectors for greatCirclePoints to reduce allocation pressure
const _startVec = new THREE.Vector3();
const _endVec = new THREE.Vector3();
const _interpVec = new THREE.Vector3();

/**
 * Generate great-circle arc points between two positions on the globe.
 * Returns elevated Vector3 arc points.
 * Optimized: reuses temporary vectors for interpolation.
 */
export function greatCirclePoints(
  startLat: number,
  startLon: number,
  endLat: number,
  endLon: number,
  segments: number = 64,
  radius: number = GLOBE_RADIUS
): THREE.Vector3[] {
  const sPos = latLonToVector3(startLat, startLon, radius);
  const ePos = latLonToVector3(endLat, endLon, radius);
  _startVec.set(sPos[0], sPos[1], sPos[2]);
  _endVec.set(ePos[0], ePos[1], ePos[2]);

  const distance = _startVec.distanceTo(_endVec);
  const maxArcHeight = arcHeight(distance, radius);

  // Pre-allocate result array
  const points = new Array<THREE.Vector3>(segments + 1);

  for (let i = 0; i <= segments; i++) {
    const t = i / segments;

    // Spherical interpolation for position on globe surface
    _interpVec.copy(_startVec).lerp(_endVec, t);
    _interpVec.normalize();

    // Elevation: parabolic arc above the surface
    const elevation = Math.sin(t * Math.PI) * maxArcHeight;
    _interpVec.multiplyScalar(radius + elevation);

    points[i] = _interpVec.clone();
  }

  return points;
}

/**
 * Calculate arc peak height based on distance between points.
 * Longer distances = taller arcs.
 */
export function arcHeight(
  distance: number,
  radius: number = GLOBE_RADIUS
): number {
  // Clamp distance relative to globe size
  const normalizedDist = Math.min(distance / (2 * radius), 1);
  // Arc height: 5% to 30% of radius depending on distance
  return radius * (0.05 + normalizedDist * 0.25);
}

/**
 * Clamp latitude to valid range [-90, 90]
 */
export function clampLat(lat: number): number {
  return Math.max(-90, Math.min(90, lat));
}

/**
 * Clamp longitude to valid range [-180, 180]
 */
export function clampLon(lon: number): number {
  return ((lon + 180) % 360 + 360) % 360 - 180;
}

/**
 * Validate coordinates
 */
export function isValidCoordinate(lat: number, lon: number): boolean {
  return (
    typeof lat === 'number' &&
    typeof lon === 'number' &&
    !isNaN(lat) &&
    !isNaN(lon) &&
    lat >= -90 &&
    lat <= 90 &&
    lon >= -180 &&
    lon <= 180
  );
}

import { getCountryInfo } from './countryNames';

// Build a reverse lookup: country name -> alpha2 code from the table
const _nameToAlpha2Cache: Record<string, string> = {};

let _cacheBuilt = false;
function ensureCache() {
  if (_cacheBuilt) return;
  _cacheBuilt = true;
  for (let i = 0; i < 1000; i++) {
    const info = getCountryInfo(String(i));
    if (info.alpha2 !== '??') {
      _nameToAlpha2Cache[info.name.toLowerCase()] = info.alpha2;
    }
  }
  // Common aliases
  _nameToAlpha2Cache['united states'] = 'US';
  _nameToAlpha2Cache['usa'] = 'US';
  _nameToAlpha2Cache['uk'] = 'GB';
  _nameToAlpha2Cache['czech republic'] = 'CZ';
  _nameToAlpha2Cache['ivory coast'] = 'CI';
  _nameToAlpha2Cache['burma'] = 'MM';
  _nameToAlpha2Cache['swaziland'] = 'SZ';
  _nameToAlpha2Cache['macedonia'] = 'MK';
}

export function getIsoCode(countryName: string): string {
  if (!countryName || countryName.startsWith('Region')) return '??';
  ensureCache();
  
  const lower = countryName.toLowerCase();
  if (_nameToAlpha2Cache[lower]) return _nameToAlpha2Cache[lower];
  
  for (const [name, code] of Object.entries(_nameToAlpha2Cache)) {
    if (lower.includes(name) || name.includes(lower)) return code;
  }

  return '??';
}
\end{verbatim}

---

\section*{Appendix H: GLSL Custom Shaders for Globe and Maps}
To achieve high performance visual rendering of geographic elements and glowing indicators on the client's screen, the Predictra Threat Map shifts processing of procedural visual elements to the GPU. This is accomplished using custom WebGL shader materials implemented inside Three.js and React Three Fiber. This appendix details the complete GLSL source codes of the custom vertex and fragment shaders.

\subsection*{H.1 Atmospheric Glow Shader}
This shader generates a smooth outer Fresnel glow surrounding the 3D procedural sphere. The vertex shader calculates light-intensity vectors relative to the user's camera view angle, and the fragment shader applies an additive glowing alpha overlay:
\begin{verbatim}
// Atmospheric Glow - Vertex Shader
varying float vIntensity;
void main() {
  vec3 vNormal = normalize(normalMatrix * normal);
  vec4 mvPosition = modelViewMatrix * vec4(position, 1.0);
  vec3 vNormel = normalize(-mvPosition.xyz);
  vIntensity = pow(0.7 - dot(vNormal, vNormel), 3.0);
  gl_Position = projectionMatrix * mvPosition;
}

// Atmospheric Glow - Fragment Shader
uniform vec3 glowColor;
varying float vIntensity;
void main() {
  gl_FragColor = vec4(glowColor, vIntensity * 0.6);
  if (gl_FragColor.a < 0.01) discard;
}
\end{verbatim}

\subsection*{H.2 Country Dot Grid Sweep Shader}
This shader renders the glowing points on geographic landmasses. It includes a time-dependent linear coordinate sweep function, blending between a primary cyan color and a secondary green accent to represent real-time security scanning sweeps across coordinates:
\begin{verbatim}
// Country Dot Grid - Vertex Shader
varying vec3 vPosition;
void main() {
  vPosition = position;
  vec4 mvPosition = modelViewMatrix * vec4(position, 1.0);
  gl_PointSize = (10.0 / -mvPosition.z) * 1.5;
  gl_Position = projectionMatrix * mvPosition;
}

// Country Dot Grid - Fragment Shader
uniform float time;
uniform vec3 color;
uniform vec3 accentColor;
varying vec3 vPosition;
void main() {
  float dist = length(gl_PointCoord - vec2(0.5));
  if (dist > 0.5) discard;
  float sweep = fract(vPosition.x * 0.2 - time * 0.15);
  float sweepIntensity = smoothstep(0.95, 1.0, sweep);
  vec3 finalColor = mix(color, accentColor, sweepIntensity * 0.8);
  gl_FragColor = vec4(finalColor, 0.2 + sweepIntensity * 0.4);
}
\end{verbatim}

\subsection*{H.3 Flat 2D Enterprise Map Grid Shader}
Used exclusively during flat map projections, this shader renders a subtle slate backgrounds with a procedural coordinate tracking grid and glowing border lines:
\begin{verbatim}
// 2D Map - Vertex Shader
varying vec2 vUv;
void main() {
  vUv = uv;
  gl_Position = projectionMatrix * modelViewMatrix * vec4(position, 1.0);
}

// 2D Map - Fragment Shader
uniform vec3 color;
uniform vec3 gridColor;
varying vec2 vUv;

void main() {
  // Subtle Procedural Coordinate Grid lines
  vec2 gUv = vUv * vec2(60.0, 30.0);
  vec2 gridLine = abs(fract(gUv - 0.5) - 0.5) / fwidth(gUv);
  float grid = 1.0 - min(min(gridLine.x, gridLine.y), 1.0);
  
  vec3 finalColor = color + gridColor * grid * 0.02; 
  
  float borderX = smoothstep(0.01, 0.0, vUv.x) + 
                  smoothstep(0.99, 1.0, vUv.x);
  float borderY = smoothstep(0.01, 0.0, vUv.y) + 
                  smoothstep(0.99, 1.0, vUv.y);
  float border = borderX + borderY;
  finalColor += gridColor * border * 0.1;

  float alpha = max(grid * 0.15, border * 0.4);
  gl_FragColor = vec4(finalColor, alpha);
}
\end{verbatim}

